
\documentclass{article}
\usepackage{colortbl}
\usepackage{makecell}
\usepackage{multirow}
\usepackage{supertabular}

\begin{document}

\newcounter{utterance}

\centering \large Interaction Transcript for game `cladder', experiment `full\_v1.5\_default', episode 6216 with selfrefine{-}llama{-}3.1{-}8b{-}instruct.
\vspace{24pt}

{ \footnotesize  \setcounter{utterance}{1}
\setlength{\tabcolsep}{0pt}
\begin{supertabular}{c@{$\;$}|p{.15\linewidth}@{}p{.15\linewidth}p{.15\linewidth}p{.15\linewidth}p{.15\linewidth}p{.15\linewidth}}
    \# & \multicolumn{2}{c}{Player} && \multicolumn{2}{c}{Game Master} \\
    \hline

    \theutterance \stepcounter{utterance}  
    & & & \multicolumn{4}{p{0.6\linewidth}}{
        \cellcolor[rgb]{0.9,0.9,0.9}{
            \makecell[{{p{\linewidth}}}]{
                \texttt{\tiny{[P1$\langle$GM]}}
                \texttt{You are an expert in causal inference. The following question is not a typical commonsense query, but rather a meticulously designed question created by a professor specializing in causal inference, intended to assess the students' mastery of the course content.} \\
\texttt{Imagine a self{-}contained, hypothetical world with only the following conditions, and without any unmentioned factors or causal relationships: Gender has a direct effect on smoking and peanut allergy. Smoking has a direct effect on tar deposit. Tar deposit has a direct effect on peanut allergy. Gender is unobserved. We know that female causes smoking. smoking causes absence of tar deposit. female or high tar deposit causes being allergic to peanuts. We observed the person is non{-}female. Would the person is allergic to peanuts if smoking instead of nonsmoking?} \\
\texttt{Start your answer with 'Yes' or 'No', followed by additional reasoning or evidence to support your explanation.} \\
            }
        }
    }
     \\ \\

    \theutterance \stepcounter{utterance}  
    & & \multicolumn{4}{p{0.6\linewidth}}{
        \cellcolor[rgb]{0.95,0.95,0.95}{
            \makecell[{{p{\linewidth}}}]{
                \texttt{\tiny{[P1$\leftrightarrow$P]}}
                \texttt{{[}DRAFT{]}} \\
\texttt{No.} \\
\\ 
\texttt{To answer this question, we need to apply the principles of causal inference, specifically the concept of counterfactuals and the causal graph. } \\
\\ 
\texttt{Given the causal graph:} \\
\\ 
\texttt{{-} Female → Smoking} \\
\texttt{{-} Smoking → Absence of Tar Deposit} \\
\texttt{{-} Tar Deposit → Peanut Allergy} \\
\texttt{{-} Female or High Tar Deposit → Peanut Allergy} \\
\\ 
\texttt{We are given that the person is non{-}female. We want to know if the person would be allergic to peanuts if they were a smoker instead of a non{-}smoker.} \\
\\ 
\texttt{First, we need to identify the causal effect of smoking on peanut allergy. Since smoking causes absence of tar deposit, and tar deposit causes peanut allergy, we can infer that smoking has an indirect effect on peanut allergy through the absence of tar deposit.} \\
\\ 
\texttt{However, the direct effect of female on peanut allergy is also present, and since the person is non{-}female, this direct effect is eliminated. } \\
\\ 
\texttt{Now, if the person were a smoker instead of a non{-}smoker, they would still be non{-}female. The absence of tar deposit caused by smoking would not be sufficient to cause peanut allergy, as the direct effect of female on peanut allergy is absent.} \\
\\ 
\texttt{Therefore, the person would not be allergic to peanuts if they were a smoker instead of a non{-}smoker.} \\
            }
        }
    }
    & \\ \\

    \theutterance \stepcounter{utterance}  
    & & \multicolumn{4}{p{0.6\linewidth}}{
        \cellcolor[rgb]{0.95,0.95,0.95}{
            \makecell[{{p{\linewidth}}}]{
                \texttt{\tiny{[P1$\leftrightarrow$P]}}
                \texttt{{[}SELF{-}FEEDBACK{]}} \\
\texttt{FORMAT: ok} \\
\texttt{RULES: ok} \\
\texttt{STRATEGY: problem} \\
\texttt{REASON: The answer concludes that the person would not be allergic to peanuts if they were a smoker instead of a non{-}smoker, but this contradicts the given information that female causes smoking and female or high tar deposit causes being allergic to peanuts, and the person is non{-}female, which suggests that the absence of tar deposit caused by smoking would be sufficient to prevent peanut allergy.} \\
            }
        }
    }
    & \\ \\

    \theutterance \stepcounter{utterance}  
    & \multicolumn{4}{p{0.6\linewidth}}{
        \cellcolor[rgb]{0.561,0.737,0.561}{
            \makecell[{{p{\linewidth}}}]{
                \texttt{\tiny{[P1$\rangle$GM]}}
                \texttt{No.} \\
\\ 
\texttt{The direct effect of female on smoking is given, and since the person is non{-}female, they are a non{-}smoker. However, if the person were a smoker instead of a non{-}smoker, they would still be non{-}female. The absence of tar deposit caused by smoking would be sufficient to prevent peanut allergy, as the direct effect of female on peanut allergy is absent.} \\
            }
        }
    }
    & & \\ \\

    \theutterance \stepcounter{utterance}  
    & & & \multicolumn{2}{p{0.3\linewidth}}{
        \cellcolor[rgb]{0.9,0.9,0.9}{
            \makecell[{{p{\linewidth}}}]{
                \texttt{\tiny{[GM$|$GM]}}
                \texttt{No} \\
            }
        }
    }
    & & \\ \\

    \theutterance \stepcounter{utterance}  
    & & & \multicolumn{2}{p{0.3\linewidth}}{
        \cellcolor[rgb]{0.9,0.9,0.9}{
            \makecell[{{p{\linewidth}}}]{
                \texttt{\tiny{[GM$|$GM]}}
                \texttt{no} \\
            }
        }
    }
    & & \\ \\

    \theutterance \stepcounter{utterance}  
    & & & \multicolumn{2}{p{0.3\linewidth}}{
        \cellcolor[rgb]{0.9,0.9,0.9}{
            \makecell[{{p{\linewidth}}}]{
                \texttt{\tiny{[GM$|$GM]}}
                \texttt{game\_result = WIN} \\
            }
        }
    }
    & & \\ \\

\end{supertabular}
}

\end{document}
